\documentclass[11pt]{article}

\usepackage{amsmath,amssymb,amsthm,amsfonts}
\usepackage{fullpage}
\usepackage{hyperref}
\usepackage{bm}
\usepackage{mathtools}
\usepackage{graphicx}
\usepackage{enumerate}
\usepackage{listings}
\usepackage{xcolor}

\lstset{
  basicstyle=\ttfamily\small,
  keywordstyle=\color{blue},
  commentstyle=\color{gray},
  stringstyle=\color{red},
  showstringspaces=false,
  breaklines=true,
  frame=single
}

\title{Multiplicity Knot Theory v3.0:\\
Prime-Weighted Braid Invariants and a Cryptographic Commitment Prototype}

\author{Helix AI Innovations\\ \& \\ Citizen Gardens \\ The Foundation of Multiplicity}

\date{\today}

\begin{document}

\maketitle

\begin{abstract}
Multiplicity Theory aims to construct braid invariants for knots and links whose structure is driven by prime-number arithmetic rather than by arbitrary quantum-group parameter choices. Building on a prior v2.1 note introducing a prime-weighted protection factor $P(K)$ for ``constitutional lattices'' \cite{MTv21}, we present a v3.0 reframing that:

\begin{itemize}
  \item defines a prime-colored braid category and a strand-dependent $R$-matrix $R_{p,q}$ derived from standard $U_q(\mathfrak{sl}_2)$ data,
  \item introduces a prime-weighted functional $Z(K)$ via a modified trace against strand-local weights $O_p O_p^\dagger$,
  \item constructs a protection functional $P(K;X)$ based on truncated prime averages $c_0(X)$ and $z(X)$ treated as conjectural limits, and
  \item sketches a multiplicity-based commitment (MBC) scheme whose binding property is heuristically tied to the injectivity and stability of the invariant pipeline.
\end{itemize}

We give precise conditional statements, explicitly record the necessary conjectures, and provide a small-scale Python prototype suitable for experimental validation.
\end{abstract}
 \newpage
\tableofcontents
\newpage

\section{Executive Summary}

The original v2.1 document claimed a parameter-free topological invariant $P(K)$, built by conjugating a standard $U_q(\mathfrak{sl}_2)$ $R$-matrix by prime-indexed operators $O_p$ and combining the resulting braid invariant $Z(K)$ with number-theoretically derived constants $c_0 = \ln 10$ and $z = 1/(2\cos 1)$ \cite{MTv21}. A subsequent self-critique identified several serious issues: an ill-defined modified Markov trace, incomplete or flawed derivations of the constants, and a lack of a canonical prime assignment to strands.

In this v3.0 program document we:

\begin{enumerate}[(i)]
  \item Restrict attention to a \emph{prime-colored braid category}, with a strand-dependent $R$-matrix
  \[
    R_{p,q} := (O_p \otimes O_q) R_{\mathrm{std}} (O_p^\dagger \otimes O_q^\dagger)
  \]
  built from standard $U_q(\mathfrak{sl}_2)$ data at $q = e^{i\pi/3}$ and prime-indexed unitaries $O_p \in SU(2)$ \cite[Sec.~2.1--2.3]{MTv21}.

  \item Define a prime-weighted functional
  \[
    Z(K) := \operatorname{Tr}\big(\rho(\beta_K) W_K\big), \qquad
    W_K = \bigotimes_{j=1}^n (O_{p_j} O_{p_j}^\dagger),
  \]
  where $\beta_K$ is a canonical prime-colored braid representative of $K$ and $p_j$ is the prime label of strand $j$.

  \item Introduce truncated prime-averaged constants
  \[
    c_0(X) = \frac{1}{2\pi} \frac{\sum_{p\le X} (p\ln p)^{-1}}{\sum_{p\le X} p^{-2}}, \qquad
    z(X) = \frac{\sum_{p\le X} 2\sin^2(\ln p)/(p\ln p)}{\sum_{p\le X} (p\ln p)^{-1}},
  \]
  and define a finite-$X$ protection functional
  \[
    P(K;X) := \exp\big(c_0(X) c(K)\big)\, |Z(K)|^{z(X)}.
  \]
  The large-$X$ limits $c_0$ and $z$ are treated as conjectural, supported by numerics rather than proved identities \cite[Sec.~3--4]{MTv21}.

  \item Formulate a two-phase multiplicity-based commitment (MBC) scheme and a clear security model, in which binding is heuristic and tied to three assumption clusters (algebraic invariants, encoding/canonicalization, and hash/invariant collisions).

  \item Provide a concrete 4-strand prototype using primes $(2,3,5,7)$ and Python/NumPy code to compute $Z(K)$ and $P(K;X)$ and to run small-scale collision and stability experiments.
\end{enumerate}

The result is not yet a full mathematical theory, but it is a \emph{coherent, testable} research program that turns the previous overclaims into explicit conjectures and experimentally accessible questions.

\section{Prime-Colored Braid Category and Local Data}

\subsection{Prime-colored objects and morphisms}

Fix an infinite ordered list of primes $(p_1,p_2,p_3,\dots)$. Objects of our category are finite ordered tuples
\[
  \mathbf{p} = (p_{i_1},\dots,p_{i_n}),
\]
which we interpret as colorings of $n$ strands by primes. Morphisms are braids on $n$ strands, with the $j$-th strand colored by $p_{i_j}$. Composition is braid concatenation; tensor product is given by juxtaposing colored sequences and braids.

\subsection{Prime-indexed local unitaries $O_p$}

For each prime $p$ we define a unitary $O_p \in SU(2)$ by
\begin{equation}
  O_p = \exp\big(i \ln p \, (\hat n_p \cdot \bm{\sigma})\big),
  \label{eq:Op-def}
\end{equation}
where $\bm{\sigma} = (\sigma_x,\sigma_y,\sigma_z)$ are the Pauli matrices and
\begin{equation}
  \hat n_p = \frac{1}{N_p}
  \begin{pmatrix}
    \sin(\ln p) \\[4pt]
    \cos(\ln p) \\[4pt]
    p^{-1/2}
  \end{pmatrix},
  \qquad
  N_p = \sqrt{1 + p^{-1}}.
  \label{eq:n-hat}
\end{equation}
The choice of $\hat n_p$ is a modeling assumption (unit vector with nontrivial dependence on $\ln p$ and $p^{-1/2}$); no deeper number-theoretic significance is currently claimed beyond this structure. A related construction appears in \cite[Eq.~(1)--(2)]{MTv21}.

\subsection{Strand-dependent $R$-matrix $R_{p,q}$}

Let $R_{\mathrm{std}}$ denote the standard $U_q(\mathfrak{sl}_2)$ $R$-matrix at $q = e^{i\pi/3}$:
\begin{equation}
  R_{\mathrm{std}} =
  \begin{pmatrix}
    q^{1/2} & 0 & 0 & 0 \\
    0 & q^{1/2} & -q^{-1/2} & 0 \\
    0 & 1 & 0 & 0 \\
    0 & 0 & 0 & q^{1/2}
  \end{pmatrix},
  \qquad q = e^{i\pi/3},
  \label{eq:R-std}
\end{equation}
as in \cite[Eq.~(4)]{MTv21}.

For a crossing of two strands colored by primes $p$ and $q$ we define
\begin{equation}
  R_{p,q} := (O_p \otimes O_q)\, R_{\mathrm{std}}\, (O_p^\dagger \otimes O_q^\dagger).
  \label{eq:R-pq}
\end{equation}
Given a prime-colored braid $\beta$ on $n$ strands with color sequence $(p_{i_1},\dots,p_{i_n})$, we obtain a representation
\[
  \rho(\beta) \in \mathrm{End}\big((\mathbb{C}^2)^{\otimes n}\big)
\]
by assigning $R_{p_{i_j},p_{i_{j+1}}}$ (or its inverse) to each elementary crossing of strands $j$ and $j+1$. This matches the strand-dependent construction sketched in \cite[Sec.~2.2]{MTv21}, but now with explicit tensor indices and fixed prime colors.

\section{Yang--Baxter Structure and Restricted Markov Invariance}

\subsection{Projective YBE conjecture}

The standard matrix $R_{\mathrm{std}}$ satisfies the Yang--Baxter equation (YBE). After conjugation by $O_p \otimes O_q$, the resulting $R_{p,q}$ is no longer guaranteed to satisfy YBE in a straightforward way, because different strands carry different $O_p$ factors \cite[Sec.~2.3]{MTv21}. We therefore formulate a projective YBE condition.

For any triple of primes $(p,q,r)$ define
\[
  L_{p,q,r} := R_{p,q} R_{p,r} R_{q,r}, \qquad
  R_{p,q,r} := R_{q,r} R_{p,r} R_{p,q}.
\]

\medskip

\noindent\textbf{Conjecture 2.1 (Projective prime-labeled YBE).}
\emph{For all primes $p,q,r$ there exists a phase $\lambda_{p,q,r} \in U(1)$ such that}
\[
  L_{p,q,r} = \lambda_{p,q,r}\, R_{p,q,r}.
\]

Empirically, the v2.1 note reports YBE residuals below $10^{-15}$ for triples in a modest prime range \cite[Sec.~2.3, Remark~2.1]{MTv21}, but no proof is given. In this v3.0 framework, Conjecture~2.1 is taken as a central algebraic assumption.

Under Conjecture~2.1, $\rho$ defines a projective braid group representation in the prime-colored category.

\subsection{Prime-weighted functional $Z(K)$}

Let $\beta_K$ be a fixed prime-colored braid representative of a knot or link $K$ on $n$ strands, with color sequence $(p_{i_1},\dots,p_{i_n})$. We define the strand-weight operator
\begin{equation}
  W_K := \bigotimes_{j=1}^n \big(O_{p_{i_j}} O_{p_{i_j}}^\dagger\big).
  \label{eq:W-K}
\end{equation}
The \emph{prime-weighted functional} is then
\begin{equation}
  Z(K) := \operatorname{Tr}\big(\rho(\beta_K) W_K\big).
  \label{eq:Z-K}
\end{equation}
This makes precise the ``modified Markov trace'' heuristic of \cite[Eq.~(6)]{MTv21}, where a product of $O_p O_p^\dagger$ factors was written but not fully tensorialized.

We do not claim that $Z$ is a Markov trace in the sense of Ocneanu; instead we seek invariance under a restricted set of moves compatible with a canonical prime assignment.

\subsection{Canonical prime assignment and restricted moves}

To ensure well-definedness we fix:

\begin{itemize}
  \item A deterministic procedure $\mathsf{Canon}$ that maps an initial diagram or braid for $K$ to a braid word $\beta_K$ on $n$ strands in a canonical form.\footnote{For the 4-strand prototype in Section~\ref{sec:prototype}, we take $\mathsf{Canon}$ to be the identity and work directly with the encoding.}
  \item A prime assignment rule: strand $j$ of $\beta_K$ is labeled by prime $p_j$ for $j=1,\dots,n$.
\end{itemize}

Allowed moves are:

\begin{itemize}
  \item Colored conjugations that do not permute strand endpoints (indices $1,\dots,n$ are fixed), so the prime labels remain aligned and cyclicity of trace preserves $Z(K)$.
  \item Terminal stabilization/destabilization that adds/removes a last strand with the next unused prime $p_{n+1}$, with normalization $\operatorname{Tr}(O_{p_{n+1}} O_{p_{n+1}}^\dagger) = 2$ to preserve $Z(K)$ under paired moves, echoing the sketch of \cite[Thm.~2.1]{MTv21}.
\end{itemize}

\medskip

\noindent\textbf{Assumption A2 (Restricted Markov invariance).}
\emph{For any two diagrams of the same knot or link $K$, the canonicalization and prime-assignment pipeline yields braids whose associated $Z(K)$ are equal up to numerical tolerance.}

\section{Prime-Harmonic Couplings and the Constant $c_0$}

\subsection{Truncated prime averages}

Let $\mu_p := 1/(p \ln p)$ be the usual ``prime-counting'' weight \cite[Sec.~1]{MTv21}. For a finite cutoff $X>0$ define the normalized truncated average
\[
  \langle f(p) \rangle_X := \frac{\sum_{p\le X} f(p)\,\mu_p}{\sum_{p\le X} \mu_p}.
\]
In v2.1, a self-consistency condition for a prime-harmonic coupling $L_p = c_0 \ln(p)/p$ and a frequency factor $\omega_p = 2\pi/\ln p$ was claimed to fix $c_0 = \ln 10$, but the derivation as written is mathematically inconsistent when extended to all primes \cite[Sec.~3.3]{MTv21}. We therefore work with a finite-$X$ definition.

\subsection{Finite-$X$ self-consistency for $c_0(X)$}

For each cutoff $X$ define $c_0(X)$ by
\begin{equation}
  \big\langle L_p \omega_p \big\rangle_X = 1, \qquad
  L_p = c_0(X) \frac{\ln p}{p},\ \omega_p = \frac{2\pi}{\ln p}.
\end{equation}
This simplifies to
\[
  c_0(X)\, 2\pi \, \frac{\sum_{p\le X} p^{-2}}{\sum_{p\le X} (p\ln p)^{-1}} = 1,
\]
so
\begin{equation}
  c_0(X) = \frac{1}{2\pi} \frac{\sum_{p\le X} (p\ln p)^{-1}}{\sum_{p\le X} p^{-2}}.
  \label{eq:c0X}
\end{equation}

\medskip

\noindent\textbf{Conjecture 3.1 (Prime-harmonic coupling limit).}
\emph{The limit $c_0 = \lim_{X\to\infty} c_0(X)$ exists and equals $\ln 10$.}

In v2.1, a numerical fit $c_0 \approx \ln 10$ with small residuals is reported \cite[Sec.~3.3]{MTv21}; here we recast that observation as a conjectural limiting value.

\section{Prime-Distribution Exponent $z$}

\subsection{Truncated prime-average for $z(X)$}

Using the same weights $\mu_p$, define
\begin{equation}
  z(X) := \left\langle 2\sin^2(\ln p) \right\rangle_X
   = \frac{\sum_{p\le X} 2\sin^2(\ln p)/(p\ln p)}{\sum_{p\le X} (p\ln p)^{-1}}.
  \label{eq:zX}
\end{equation}
The v2.1 note numerically observes $z \approx 0.9248$ for $p\le 10^6$, matching $1/(2\cos 1) \approx 0.9254$ to within about $0.06\%$ \cite[Sec.~4]{MTv21}.

\medskip

\noindent\textbf{Conjecture 4.1 (Prime-zeta exponent).}
\emph{The limit $z := \lim_{X\to\infty} z(X)$ exists and equals $1/(2\cos 1)$.}

Again, this is explicitly conjectural: the exact evaluation of the associated prime-zeta-like integral remains open \cite[Remark~4.1]{MTv21}.

\section{Protection Functional and Conditional Main Result}

\subsection{Finite-$X$ protection functional}

Given a knot or link $K$ with canonical prime-colored braid $\beta_K$ and corresponding $Z(K)$, define the finite-$X$ protection functional
\begin{equation}
  P(K;X) := \exp\big(c_0(X) c(K)\big)\, |Z(K)|^{z(X)},
  \label{eq:PKX}
\end{equation}
where $c(K)$ is the crossing number of a chosen canonical diagram of $K$. This mirrors the structure of \cite[Eq.~(7)]{MTv21}, but with $c_0$ and $z$ now explicitly cutoff-dependent and conjectural in the limit.

\subsection{Conditional prime-weighted invariant}

We can now state a main conditional theorem, in the style of Theorem 1.1 from our earlier draft.

\medskip

\noindent\textbf{Theorem 6.1 (Conditional prime-weighted protection invariant).}
\emph{Fix:}

\begin{itemize}
  \item the ordered prime list $(p_1,p_2,\dots)$,
  \item the data $(O_p)$ and $(R_{p,q})$ as in \eqref{eq:Op-def}--\eqref{eq:R-pq},
  \item a deterministic canonicalization procedure $\mathsf{Canon}$ and prime assignment rule as in Section~3.2,
  \item a finite set of cutoffs $\{X_1,\dots,X_t\}$.
\end{itemize}

\emph{Assume:}
\begin{enumerate}[(C1)]
  \item \textbf{(Projective YBE)} Conjecture~2.1 holds for all primes $p,q,r$.
  \item \textbf{(Restricted invariance)} $Z(K)$ is invariant (up to numerical tolerance) under the moves used by $\mathsf{Canon}$ on equivalent diagrams of $K$ (Assumption~A2).
  \item \textbf{(Convergence)} The limits $c_0 = \lim_{X\to\infty} c_0(X)$, $z = \lim_{X\to\infty} z(X)$ exist and are finite (Conjectures~3.1--4.1).
\end{enumerate}

\emph{Then:}
\begin{enumerate}[(1)]
  \item For each finite $X_i$, the quantity $P(K;X_i)$ in \eqref{eq:PKX} is a well-defined numerical invariant of $K$ under isotopy, relative to the fixed prime list and canonicalization.
  \item The limit $P(K) = \lim_{X\to\infty} P(K;X)$ exists and defines a conjectural prime-weighted protection invariant of $K$.
\end{enumerate}

\section{Multiplicity-Based Commitment (MBC) Scheme}

We now describe a non-interacting commitment scheme $\mathsf{MBC} = (\mathsf{Setup},\mathsf{Commit},\mathsf{Open})$ built on top of $P(K;X)$. The goal is not to claim production-ready security but to align \emph{cryptographic binding} with the injectivity and stability of the v3.0 invariant pipeline.

\subsection{Syntax}

\begin{itemize}
  \item $\mathsf{Setup}(1^\lambda) \to \mathsf{pp}$: generates public parameters $\mathsf{pp}$.
  \item $\mathsf{Commit}(\mathsf{pp}, m) \to (C,\mathsf{aux})$: on input message $m$, outputs a commitment $C$ and opening data $\mathsf{aux}$.
  \item $\mathsf{Open}(\mathsf{pp}, C, m, \mathsf{aux}) \to \{0,1\}$: verifies a purported opening and returns $1$ for accept, $0$ otherwise.
\end{itemize}

\subsection{Setup}

$\mathsf{Setup}(1^\lambda)$:

\begin{enumerate}
  \item Fix a global prime list $(p_j)_j$, and the operators $O_p$, $R_{p,q}$, and cutoffs $\{X_i\}$ as in previous sections.
  \item Fix an encoding $\mathrm{Enc}: \{0,1\}^\ell \to \{\text{initial diagrams/braids}\}$ and canonicalization $\mathsf{Canon}$.
  \item Fix a collision-resistant hash function $H$.
\end{enumerate}

Set
\[
  \mathsf{pp} := \big( (p_j)_j, (O_p)_p, (R_{p,q})_{p,q}, \{X_i\}_i, H, \mathrm{Enc}, \mathsf{Canon} \big).
\]

\subsection{Commit}

$\mathsf{Commit}(\mathsf{pp}, m)$:

\begin{enumerate}
  \item Compute $D_m := \mathrm{Enc}(m)$ and $(\beta_m,n_m) := \mathsf{Canon}(D_m)$.
  \item Assign primes $p_1,\dots,p_{n_m}$ to strands, construct $Z(K_m)$ and $P(K_m;X_i)$ for each cutoff.
  \item Serialize $\beta_m$ as $\sigma_m$ and set $T_m := H(\sigma_m \,\|\, m)$.
\end{enumerate}

Output
\[
  C_m := \big(T_m,\{P(K_m;X_i)\}_i\big), \quad \mathsf{aux}_m := \sigma_m.
\]

\subsection{Open}

$\mathsf{Open}(\mathsf{pp}, C, m, \mathsf{aux})$ recomputes the canonical braid from $m$, verifies $T_m$ and the invariant vector $\{P(K_m;X_i)\}$ against the published $C$, and returns $1$ if all checks pass within tolerance, $0$ otherwise.

\subsection{Binding game}

We formalize the binding game $\mathsf{BindGame}_{\mathsf{MBC}}(\mathcal{A},\lambda)$ where an adversary $\mathcal{A}$ gets $\mathsf{pp}$, outputs $(C,m,\mathsf{aux},m',\mathsf{aux}')$, and wins if $m\ne m'$ and both openings verify. The binding advantage is
\[
  \mathbf{Adv}^{\mathsf{bind}}_{\mathsf{MBC}}(\mathcal{A},\lambda) = \Pr\big[ \mathsf{BindGame}_{\mathsf{MBC}}(\mathcal{A},\lambda)=1\big].
\]

Under assumptions about hash collision-resistance, injectivity of the encoding/canonicalization map $m\mapsto \sigma_m$, and the absence of ``cheap'' invariant collisions $P(K_m;X_i) = P(K_{m'};X_i)$ for $m\ne m'$, any double opening would either break the hash function or witness a nontrivial collision in the invariant pipeline (or a failure of the invariant's stability), making the binding advantage heuristically negligible.

A more detailed security model and assumption set (Assumptions A--C) can be stated following the structure in the chat-derived draft.

\section{Prototype Implementation and Experimental Plan}
\label{sec:prototype}

We now describe a concrete 4-strand prototype with primes $(2,3,5,7)$ and a 12-bit message encoding, and present a Python-style pseudocode snippet that instantiates the v3.0 constructions. This mirrors the structure and constants of the v2.1 note but in a controlled, small-parameter setting \cite[Table~2]{MTv21}.

\subsection{Prototype parameters}

\begin{itemize}
  \item Strand count: $n=4$ with primes $(2,3,5,7)$.
  \item Base $R_{\mathrm{std}}$ and $O_p$ as in \eqref{eq:R-std}--\eqref{eq:Op-def}.
  \item Cutoffs: e.g.\ $X\in\{10^3,10^4\}$ for $c_0(X)$ and $z(X)$.
  \item Encoding $\mathrm{Enc}$: 12-bit messages split into four 3-bit blocks; each block $m^{(j)}$ selects a small braid gadget on 4 strands and an exponent. The full braid is the concatenation of four such gadgets.
\end{itemize}

\subsection{Representative code snippet}

The following excerpt shows the core of the prototype: construction of $O_p$, $R_{p,q}$, the braid representation $\rho(\beta)$ on 4 strands, and the computation of $Z(K)$ and $P(K;X)$.

\begin{lstlisting}[language=Python,caption={Core v3.0 prototype code (4-strand case)}]
import numpy as np
import cmath
from math import log, sqrt, sin, cos
import sympy as sp

DTYPE = np.complex128

def kron(*matrices):
    result = matrices[0]
    for M in matrices[1:]:
        result = np.kron(result, M)
    return result

# Pauli matrices
sigma_x = np.array([[0, 1], [1, 0]], dtype=DTYPE)
sigma_y = np.array([[0, -1j], [1j, 0]], dtype=DTYPE)
sigma_z = np.array([[1, 0], [0, -1]], dtype=DTYPE)

def n_hat_p(p):
    ln_p = log(p)
    v = np.array([sin(ln_p), cos(ln_p), p**-0.5], dtype=float)
    Np = sqrt(1.0 + 1.0/p)
    return v / Np

def O_p(p):
    ln_p = log(p)
    nx, ny, nz = n_hat_p(p)
    n_dot_sigma = nx*sigma_x + ny*sigma_y + nz*sigma_z
    theta = ln_p
    I2 = np.eye(2, dtype=DTYPE)
    return np.cos(theta)*I2 + 1j*np.sin(theta)*n_dot_sigma

def R_std():
    q = cmath.exp(1j * cmath.pi / 3.0)
    q_half = q**0.5
    q_minus_half = q**-0.5
    R = np.zeros((4,4), dtype=DTYPE)
    R[0,0] = q_half
    R[1,1] = q_half
    R[1,2] = -q_minus_half
    R[2,1] = 1.0
    R[3,3] = q_half
    return R

R_STD = R_std()
PRIMES_4 = (2, 3, 5, 7)

def R_pq(p, q):
    Op = O_p(p); Oq = O_p(q)
    U = kron(Op, Oq)
    U_dag = U.conj().T
    return U @ R_STD @ U_dag

# Precompute generators
I2 = np.eye(2, dtype=DTYPE)
GEN = {}
for j in [1,2,3]:
    p_left = PRIMES_4[j-1]
    p_right = PRIMES_4[j]
    R = R_pq(p_left, p_right)
    R_inv = np.linalg.inv(R)
    if j == 1:
        op_plus  = kron(R, I2, I2)
        op_minus = kron(R_inv, I2, I2)
    elif j == 2:
        op_plus  = kron(I2, R, I2)
        op_minus = kron(I2, R_inv, I2)
    else:
        op_plus  = kron(I2, I2, R)
        op_minus = kron(I2, I2, R_inv)
    GEN[(j, +1)] = op_plus
    GEN[(j, -1)] = op_minus

class Braid:
    def __init__(self, word):
        self.word = list(word)
    def serialize(self):
        return ";".join(f"{j}:{s}" for (j,s) in self.word)

def rho(braid):
    M = np.eye(16, dtype=DTYPE)
    for (j, s) in braid.word:
        M = GEN[(j,s)] @ M
    return M

def W_K():
    factors = []
    for p in PRIMES_4:
        Op = O_p(p)
        factors.append(Op @ Op.conj().T)
    return kron(*factors)

W_GLOBAL = W_K()

def Z_of_braid(braid):
    return np.trace(rho(braid) @ W_GLOBAL)

def c0_of_X(X):
    primes = list(sp.primerange(2, X+1))
    S1 = sum(1.0/(p**2) for p in primes)
    S2 = sum(1.0/(p*log(p)) for p in primes)
    return (1.0/(2.0*cmath.pi)) * (S2 / S1)

def z_of_X(X):
    primes = list(sp.primerange(2, X+1))
    num = sum(2.0*(sin(log(p))**2)/(p*log(p)) for p in primes)
    den = sum(1.0/(p*log(p)) for p in primes)
    return num / den

def P_of_braid_X(braid, X):
    Z = Z_of_braid(braid)
    cK = len(braid.word)  # proxy for crossing number
    c0 = c0_of_X(X)
    zX = z_of_X(X)
    return np.exp(c0 * cK) * (abs(Z) ** zX)
\end{lstlisting}

Using this prototype, one can implement a commitment function by hashing the serialized braid and the message, and binding it to the vector of $P(K;X)$ values for chosen cutoffs. A small-scale experimental plan (e.g.\ enumerating $2^{12}$ messages, checking for collisions in $(\beta,P)$ pairs, and measuring inversion difficulty) directly probes the assumptions underlying binding and the stability of the invariants.

\section{Conclusion}

The v3.0 reframing of Multiplicity Theory:

\begin{itemize}
  \item isolates a mathematically meaningful core (prime-colored braid representations and prime-weighted functionals),
  \item demotes previously overclaimed constants to conjectural limits based on truncated prime statistics,
  \item introduces a structured, falsifiable cryptographic use case (MBC) whose security hinges on the same algebraic and numerical properties,
  \item and provides a concrete numerical environment in which the central conjectures can be probed.
\end{itemize}

Future work includes: a rigorous analysis of the prime-labeled YBE, analytic control of the prime-averaged constants $c_0$ and $z$, a full proof of restricted Markov invariance for $Z(K)$, and exploration of stronger cryptographic constructions leveraging the same multiplicity-theoretic structure.

\nocite{*}
\bibliographystyle{plainnat}
\bibliography{references}

\end{document}